\section{Introduction}
\label{sec:intro}

Security claims about machine learning appear faster than the field can
compare
them. Arp et al.\ read \DDpapers{} papers from CCS, IEEE S\&P, USENIX Security
and NDSS across a decade, name \DDpitfalls{} recurring pitfalls, and find
every one of the \DDpapers{} papers affected by at least \DDminPitfalls{} of
them
\cite{AQPWPWCR22}. One pitfall is the threat model. They measure practice, and
the same defect appears in the results themselves wherever several groups have
studied one attack.

Membership inference makes the defect concrete, and we use it in
Section~\ref{sec:problem} to state three problems. The attack decides whether
one record entered a model's training set, and Meeus et al.\ count at least
\MIAnewMethods{} new methods for it proposed in 2024 alone \cite{MSJFRM25}. The
published verdicts span the full range. Carlini et al.\ report that the attack
succeeds at low false positive rates \cite{CCNSTT22}. Duan et al.\ run five
attacks against models whose training set they know, and report that none
exceeds \MIAmimirAUC{} area under the receiver operating characteristic curve
(AUC) outside one domain \cite{DSMMSZTCEH24}. Das et al.\ report that on
\MIAblindSets{} published benchmarks a classifier that never queries the model
beats the best attack that does \cite{DZT25}. Zhang et al.\ argue that the
attack cannot establish training-set membership for a deployed model at all
\cite{ZDKT25}. Each result holds on its own terms. They differ in what they
let the attacker know and spend, and none of them states that in a form the
others can be checked against.

Membership inference is one attack among many. It appears in \NumMIApages{} of
the \NumSources{} works we survey and is the attack studied in
\NumMIAfamily{}. The corpus also covers evasion, poisoning, backdoors, data
reconstruction and model extraction, and on the other side computational
differential privacy, statistical-query and low-degree hardness,
bounded-memory adversaries and fine-grained cryptography. No one attack runs
through every part of the survey, so Section~\ref{sec:findings-tables} opens
with a single work read through every field and each finding names its own
cases.
% claim: membership-inference-verdicts-disagree

The field describes its adversaries and does not price them. The standing
systematization of the area sorts attacks along three parts, namely the point
in the pipeline the adversary attacks, what it is able to do once there, and
the goal it pursues \cite{PMSW18}. The middle part takes a different form in
each phase. Its Figure 3 states that adversaries act ``at inference time by
exploiting model internal information (white box) or probing the system to
infer system vulnerabilities (black box)'', and that they ``use read or write
access to the training data to mimic or corrupt the model''. None of the three
parts is a computational resource. Three later systematizations report that
their own subfield lacks a formal definition of attack success, on adversarial
examples
\cite{FK24}, on model extraction \cite{SSEP24}, and on data reconstruction
\cite{WLBZ25}. The first of the three is a preprint. Each proposes a repair, and
each repair supplies knowledge, access or a success metric. A reader holding
two of these results cannot ask which adversary was stronger, because neither
result answers it.
% claim: resource-axis-absent-from-prior-systematizations

Cryptography holds the objects that comparison needs. An adversary class fixes
what the attacker may run, and a budget fixes how much of it the attacker may
spend. A success criterion together with a tolerance settles what counts as
winning, since an attack that recovers a record approximately has recovered
it. A baseline gives the success rate available from public knowledge alone,
without which an advantage is a number with nothing to subtract. We carry
those four objects into machine learning security and read the field through
them.

We record the budget per party and not per adversary. Recording it per party
lets one statement hold a bound on the attacker, a bound on the party that
trains, a bound on the party that audits a released model, and a relation
between two of them. It also makes the most common answer visible, which is
that no party's resources enter its main claim at all.

We claim the following main contributions.

\begin{itemize}
  \item A classification of \NumSources{} formal works on the security and
privacy of learning, by the attack family together with the party each work
constrains, recording which party's resources enter its main claim, in what
unit, and whether the work treats that bound as its object
(Section~\ref{sec:findings-tables}).
% claim: classification-of-two-hundred-works
  \item An evidence layer under that classification. Every work carries the
verbatim sentence that decided it and the place in the paper it came from, or,
where a work states no bound, the places we checked. We verified each
quotation inside its own document (Section~\ref{sec:method}).
% claim: evidence-layer-per-row
  \item Six findings the classification supports, the first survey of this area
to report which resource a work bounds and which party it binds
(Section~\ref{sec:findings-tables}). We name what is dense, what is empty, and
what the emptiness means.
% claim: resource-axis-absent-from-prior-systematizations
\end{itemize}
